\def\topic{many--valued}
\ifdefined\MyHeaderIncluded
\else
\def\qwertyflag{true}

\ifdefined\MyHeaderIncluded

\else
  \newcommand{\foobar}{FooBar}
  \def\MyHeaderIncluded{true}
% if is till the end of the file

\documentclass[10pt,a4paper]{article}
\usepackage[utf8]{inputenc}
%\usepackage[german]{babel}
%\usepackage[cyrillic]{babel}
\usepackage[T1]{fontenc}
%\usepackage[T2A]{fontenc}
%\usepackage[english, russian]{babel}
\usepackage{amsmath}
\usepackage{amsfonts}
\usepackage{amssymb}
\usepackage{amsthm}
\usepackage{mathtools} %coloneqq

\usepackage{bussproofs} %

%\renewcommand{\thesubsection}{\arabic{subsection}}
%\newtheorem{thm}{Theorem}[section]
%\renewcommand{\thethm}{\arabic{subsection}.\arabic{thm}}

\usepackage{thmtools}
\usepackage{hyperref}
\hypersetup{
    colorlinks=true,
    linkcolor=blue,
    filecolor=magenta,      
    urlcolor=cyan,
    pdftitle={Overleaf Example},
    %pdfpagemode=FullScreen,
    }
%\usepackage{stix-type1}%\turnediota
\theoremstyle{definition}

\newtheorem{constructor}{Constructor}[section]
\newtheorem{iaxiom}{Axiom}
\newtheorem{axiom}{Axiom}
\newtheorem*{axiom*}{Axiom}
\newtheorem{definition}{Definition}[section]
\newtheorem*{definition*}{Definition}
\newtheorem{remark}{Remark}[section]
\newtheorem*{remark*}{Remark}
\newtheorem{recall}{Recall}[section]
\newtheorem*{recall*}{Recall}

\newtheorem{lemma}{Lemma}[section]
\newtheorem{theorem}{Theorem}[section]
\newtheorem*{theorem*}{Theorem}
\newtheorem{note}{Note}[section]
\usepackage{etoolbox}

%\newcommand*{\@Dashv}[2]{%
%  \reflectbox{$\m@th#1#2$}%
%}
\newcommand*{\separation}[3]{% 1)binder, 2)set, 3)condition
  \{#1\in #2: #3\}
}

% https://www.freecodecamp.org/news/simple-regex-tricks-for-beginners-3acb3fa257cb/
%  \\mathcal\{T\} |---> \topology
%  \\mathcal\{([^\}]*)\}   |--->  $1
\newcommand{\AND}{\ \ \&\ \ }
\newcommand{\any}{\mathrm{\ any}}
\newcommand{\borel}{\mathcal{B}}
\newcommand{\can}{\mathrm{can}}
\newcommand{\canonize}{\mathrm{canonize}}
\newcommand{\closure}{\mbox{cl}}
\newcommand{\complex}{\mathbb{C}}
\newcommand{\converges}{\longrightarrow}
\newcommand{\countable}{\mbox{\ countable}}
\newcommand{\ldefined}{\mbox{\ defined}}
\newcommand{\dom}{\mbox{dom}}
\newcommand{\Dom}{\mbox{Dom}}
\newcommand{\Done}{\noindent \textbf{Done.}\\}
\newcommand{\lexists}{\mbox{\ exists}}
\newcommand{\lforall}{\mbox{\ forall}}
\newcommand{\lunique}{\mbox{\ unique}}
\newcommand{\family}{\mathrm{\ family}}
\newcommand{\Fm}{\mathrm{Fm}}
\newcommand{\FuncSymb}{\mathrm{FuncSymb}}
\newcommand{\function}{\mathrm{\ function}}
\newcommand{\functional}{\mbox{\ functional}}
\newcommand{\IFF}{\ \Longleftrightarrow\ }
\renewcommand{\iff}{\ \longleftrightarrow\ }
\newcommand{\IMPLIES}{\ \Longrightarrow\ }
\renewcommand{\implies}{\ \longrightarrow\ }
\newcommand{\Inf}{\mathrm{Inf}}
\newcommand{\injective}{\mbox{\ injective}}
\newcommand{\injection}{\mbox{\ injection}}
\newcommand{\interior}{\mbox{int}}
\newcommand{\intersection}{\bigcap}
\newcommand{\isempty}{\mathrm{\ empty}}
\newcommand{\LB}{\mathrm{LB}}
\newcommand{\nghbset}{\mathcal{N}}
\newcommand{\nonempty}{\mathrm{\ nonempty}}
\newcommand{\Ord}{\mathrm{Ord}}
\newcommand{\power}{\mathcal{P}}
\newcommand{\pr}{\mathrm{pr}}
\newcommand{\fst}{\mathrm{fst}}
\newcommand{\snd}{\mathrm{snd}}
\newcommand{\PredSymb}{\mathrm{PredSymb}}
\newcommand{\Proof}{\noindent \textbf{Proof.}\\}
\newcommand{\ran}{\mathrm{ran}}
\newcommand{\Ran}{\mathrm{Ran}}
\newcommand{\real}{\mathbb{R}}
\newcommand{\relation}{\mathrm{\ relation}}
\newcommand{\set}{\mathrm{\ set}}
\newcommand{\singleton}{\mathrm{\ singleton}}
\newcommand{\semantics}{\mathrm{semantics}}
\newcommand{\semialgebra}{\mathcal{S}}
\newcommand{\Sup}{\mathrm{Sup}}
\newcommand{\surjective}{\mbox{\ surjective\ on\ }}
\newcommand{\Term}{\mathrm{Term}}
\newcommand{\Type}{\mathrm{Type}}
\newcommand{\topology}{\mathcal{T}}
\newcommand{\total}{\mbox{\ total\ on\ }}
\newcommand{\transitive}{\mathrm{\ transitive}}
\newcommand{\union}{\bigcup}
\newcommand{\lTrue}{\mathrm{True}}
\newcommand{\lFalse}{\mathrm{False}}
\newcommand{\lUncertain}{\mathrm{Uncertain}}
\newcommand{\lif}{\mathrm{if}}
\newcommand{\lthen}{\mathrm{then}}
\newcommand{\lelse}{\mathrm{else}}
\newcommand{\lexception}{\mathrm{exception}}
\newcommand{\lundef}{\mathrm{undef}}
\newcommand{\univocal}{\mathrm{\ univocal}}
\newcommand{\univocalat}{\mathrm{\ univocal\ at\ }}
\newcommand{\UB}{\mathrm{UB}}
\newcommand{\Un}{\mathcal{U}}
\newcommand{\Var}{\mathrm{Var}}

\newcommand{\notvarphi}{\phi}
\renewcommand{\phi}{\varphi}

\newcommand{\comment}[1]{\quad\mbox{(#1)}}
% local commands
%\preto\section{\undef\blipo} %etoolbox
\preto\subsection{\undef\blipo} %etoolbox
% https://tex.stackexchange.com/questions/105684/limit-macro-newcommand-scope-to-section
\usepackage{graphicx}
\usepackage{authblk}

\setlength\parindent{0pt} % whole doc of \noindent
\usepackage{soul}
%\usepackage{mathgl}

\fi % guard
\ifdefined\topic
\else
\def\topic{No Titel}
\fi
\title{\topic}
\author[1]{Georgy Dunaev}
\begin{document}
\maketitle
\fi
\section{Many-valued semantic experiments}
These notes explore ways to distinguish absence of a value from variation
among possible values. One proposal tracks presence and stability
separately; another interprets a predicate by the assignments on which it
holds. In the latter proposal, correlations between subformulas matter,
so a single truth table on the labels ``true'', ``false'', and ``uncertain''
need not determine compound formulas.

An open formula varying with its assignment is different from a formula
without a classical truth value under a fixed assignment. The following
small models make that distinction precise. They develop parts of the
proposal; they do not assert soundness of an unspecified descriptor
calculus. The original working notes remain below in the source.

\subsection{Positive evidence, negative evidence, and absence}
Work in the metatheory with a set $D$ of possible values and a relation
$R\subseteq D\times D$. A set $J\subseteq D\times D$ records the
\emph{jointly admissible} values of two terms. Independent candidates
$S,T\subseteq D$ give $J=S\times T$, but correlations may require a
smaller joint set. Define
\[
\begin{split}
E_R^+(J)&\iff J\neq\emptyclass\ \land\ J\subseteq R,\\
E_R^-(J)&\iff J\neq\emptyclass\ \land\ J\cap R=\emptyclass.
\end{split}
\]
These are assertions of uniform positive and uniform negative evidence.
For $J=S\times T$, the first is precisely the guarded universal relation
from the alternative-semantics section. The second uses the complement
of $R$ on the same candidate domain. It is \emph{not} the classical
negation of the first.

\begin{theorem}[Four mutually exclusive candidate states]
Exactly one of the following holds:
\begin{enumerate}
\item $J=\emptyclass$ (no jointly admissible value);
\item $E_R^+(J)$ (uniform positive evidence);
\item $E_R^-(J)$ (uniform negative evidence);
\item $J\cap R\neq\emptyclass$ and $J\setminus R\neq\emptyclass$
      (mixed evidence).
\end{enumerate}
If $\emptyclass\neq J'\subseteq J$, either uniform-evidence assertion
true of $J$ remains true of $J'$.
\end{theorem}
\Proof
The disjoint sets $J\cap R$ and $J\setminus R$ partition $J$. If both
are empty, $J$ is empty. If exactly one is nonempty, its side gives the
corresponding uniform-evidence assertion. If both are nonempty, the
state is mixed. This proves exhaustiveness and exclusivity. A subset
of $R$ remains a subset of $R$, and a subset of a set disjoint from $R$
remains disjoint from $R$. Nonemptiness of $J'$ retains the existence
guard in both cases.
\Done

For example, take $D=\{0,1\}$ and the equality relation
$R=\{(0,0),(1,1)\}$. Independent candidates for two terms, each with
candidate set $D$, give $J=D\times D$ and mixed evidence. Repeated
occurrences of a \emph{single shared value} instead give
$J=\{(0,0),(1,1)\}$ and uniform positive evidence. Thus the marginal
candidate sets do not determine the truth of equality. With
$J=\emptyclass$, both evidence assertions are false; with mixed $J$,
both are again false, for a different reason. Replacing negative evidence
by classical negation would conflate these two failures with uniform
negative evidence.

\subsection{Truth sets with an explicit absence set}
A related experiment keeps track of the assignments themselves. Let
$\Omega$ be a nonempty set of admissible assignments. Interpret each
partial term $t$ by a partial map $t:\Omega\rightharpoonup D$. For a
binary atomic formula $R(s,t)$ put
\[
\begin{split}
H_{R(s,t)}&=\dom(s)\cap\dom(t),\\
T_{R(s,t)}&=\{\omega\in H_{R(s,t)}:(s(\omega),t(\omega))\in R\},\\
F_{R(s,t)}&=H_{R(s,t)}\setminus T_{R(s,t)},\\
A_{R(s,t)}&=\Omega\setminus H_{R(s,t)}.
\end{split}
\]
Here $H$ records presence of all required values, $T$ and $F$ record
classical truth and falsity where values are present, and $A$ records
absence. These are sets in the metatheory, not additional classes in the
object language.

One possible propositional extension makes absence contagious. Define
negation by exchanging $T$ and $F$ and retaining $A$. Define conjunction
and disjunction by
\[
\begin{split}
T_{\varphi\land\psi}&=T_\varphi\cap T_\psi,\\
F_{\varphi\land\psi}&=(F_\varphi\cap H_\psi)
                         \cup(H_\varphi\cap F_\psi),\\
T_{\varphi\lor\psi}&=(T_\varphi\cap H_\psi)
                         \cup(H_\varphi\cap T_\psi),\\
F_{\varphi\lor\psi}&=F_\varphi\cap F_\psi,\\
A_{\varphi\land\psi}&=A_{\varphi\lor\psi}=A_\varphi\cup A_\psi,
\end{split}
\]
where $H_\theta=T_\theta\cup F_\theta$. This is a precise local version
of the draft's strict presence convention. The metatheory here is
classical; neither $\Omega$ nor the sets of present assignments need be
finite.

\begin{lemma}[Partition and classical restriction]
For every formula of this propositional fragment, $T_\varphi$,
$F_\varphi$, and $A_\varphi$ partition $\Omega$. On assignments where
all its atomic subformulas are present, its interpretation is classical.
In particular,
\[
T_{\varphi\lor\neg\varphi}=H_\varphi,
\qquad
A_{\varphi\lor\neg\varphi}=A_\varphi.
\]
\end{lemma}
\Proof
Atomic formulas have the partition property by definition. Negation
preserves it. For a binary connective, outside
$H_\varphi\cap H_\psi$ the assignment lies in the specified absence
set. Inside that intersection each input has exactly one classical truth
value, and the displayed clauses are the ordinary Boolean clauses.
This proves the partition property and classical restriction by induction.
For the last display, $T_\varphi$ and $F_\varphi$ exchange under negation,
so the disjunction's truth set is their union and its absence set is
unchanged.
\Done

This separates two reasons for an informal label ``uncertain''. If a
formula is present everywhere but true on only part of $\Omega$, its
excluded middle is true everywhere. If it is absent somewhere, strict
excluded middle remains absent there. Even among everywhere-present
formulas, truth-set correlations matter: on
$\Omega=\{\omega_0,\omega_1\}$ take $T_p=\{\omega_0\}$.
Both $p$ and $\neg p$ vary, yet $p\lor\neg p$ is always true whereas
$p\lor p$ still varies. A summary recording only ``present and varying''
cannot determine disjunction.

\begin{note}[Why guarding a partial value depends on the logic]
Under the strict convention just defined, an implication with a false
antecedent and an absent consequent is absent, if implication abbreviates
$\neg\varphi\lor\psi$. In PFOL, by contrast, primitive atomic predicates
with undefined arguments are false and the logical connectives remain
classical. There $f(x)\downarrow\to f(x)=f(x)$ is true even off the
domain. Thus the draft's desire for this guarded formula to hold selects
a different treatment of absence; it is not a consequence of strict
presence propagation. Neither convention changes what happens at
assignments where all values are present.
\end{note}
\begin{proposition}[Deduction under strict absence propagation]
For this propositional fragment and any set $\Gamma$ of formulas, write
$\Gamma\models_T\theta$ when
every assignment on which all formulas in $\Gamma$ are true belongs
to $T_\theta$. Put $G=\bigcap_{\gamma\in\Gamma}T_\gamma$, with
$G=\Omega$ for the empty context. If implication abbreviates
$\neg\varphi\lor\psi$, then
\[
\Gamma\models_T(\varphi\implies\psi)
\]
holds exactly when both of the following conditions hold:
\[
G\subseteq H_\varphi\cap H_\psi,
\qquad
\Gamma\cup\{\varphi\}\models_T\psi.
\]
\end{proposition}
\Proof
The strict truth clauses give
\[
T_{\varphi\implies\psi}
 =(F_\varphi\cap H_\psi)\cup(H_\varphi\cap T_\psi).
\]
In particular, this set is contained in $H_\varphi\cap H_\psi$.
If $G$ is contained in it, then an element of $G\cap T_\varphi$
cannot belong to $F_\varphi$, so it must belong to $T_\psi$.
This is the stated consequence with the extra assumption.
Conversely, suppose both conditions hold and $\omega\in G$.
Presence makes $\varphi$ either true or false at $\omega$. In the
false case, $\omega\in F_\varphi\cap H_\psi$; in the true case,
the consequence condition gives $\omega\in H_\varphi\cap T_\psi$.
Thus $\omega\in T_{\varphi\implies\psi}$.
\Done

For a concrete failure of the unqualified deduction principle, take one
assignment where $\varphi$ is false and $\psi$ is absent. Then
$\{\varphi\}\models_T\psi$ holds vacuously, but
$\models_T(\varphi\implies\psi)$ fails because the implication is
absent. Modus ponens remains sound: whenever both $\varphi$ and its
implication to $\psi$ are true, the displayed truth-set formula forces
$\psi$ to be true. Thus keeping an inference rule does not automatically
justify its usual deduction-theorem transformation in this local logic.

\begin{comment}
\section{Many--valued ideas}

To denote the constructor of an uncertain formula, we will use the symbol $\uncert$, because it is symmetrical (neither $\bot$, nor $\top$ are  symmetrical ). And it looks like crossed D, and "defined" starts with D, and it is connected with not being defined: normally(but not always!), if $f$ is a functional symbol, $f(...,\lundef,...)$.\\
\[
I_\Term(\uncert) \coloneqq U
\]

============

Remark: actually we can use 6 truth values.\\
 True = (present, stable(True))\\
False = (present, stable(False))\\
U1    = ( absent, stable(True))\\
U2    = ( absent, stable(False))\\
U3    = (present, unstable)\\
U4    = ( absent, unstable)\\
Connectives are contagious in the first argument, it means the following:\\
Let's say \# is an interpretation of a connective, and v is a truth value, then\\
 (absent, ...) \# v = (absent, ...)\\
 (note, that "..." in lhs and rhs are not necessary the same)\\
 (..., stable(True)) $\land$ (..., stable(True)) =  (..., stable(True))\\
But two predicates can stabilize each other, so these truth values does not have a truth table, that's why we will not use them.\\

Also therefore Uncertain OR Uncertain is not necessary Uncertain in such semantics!
Therefore we need to require even more: the certainity of every subformula, otherwise it will complicate the system too much?\\

e.g. $(f(x)=5) \lor (x\notin\dom(f))$.\\
Or will it require too much to require?\\
Yes, it will be too much, since then\\
$(f\function)\land(x\in\dom(f)) \implies (f(x)=f(x))$\\
will not work normally.\\
Therefore we need something more interesting like:\\
 $(\lFalse\implies\lUncertain) \ceq \lTrue$.\\
Or even brave \& better: reject the whole idea of having a defined truth table for uncertain values(or at least having a full/complete truth table).
Such idea was wrong, because it is not true even in completely classical logic of existence: the formula "$(x\in y)$" is not valid, so it has no any particular defined truth value.\\

=================


============================================\\
Idea: Truth values of predicates are subconglomerates of $(2\times)\Un\times\dots\times\Un$ of values, where $I(P)$ is true.\\
"(Always)True" corresponds to the whole subconglomerate $(1\times)\Un\times\dots\times\Un$.\\
"(Always)False" corresponds to the empty subconglomerate $(0\times)\emptyclass\times\dots\times\emptyclass$, which is simply $\emptyclass$.\\
"Uncertain" corresponds to any other subconglomerate.\\
All products here are given $n$ times: $n$ is a valency of the predicate symbol $P$.
============================================\\

For example, the following predicates:\\
\begin{itemize}
\item (of valency 0):\\
"$\bot$" corresponds to 0.("False")\\
"$\top$" corresponds to 1.("True")\\
"$\uncert$" corresponds to 2.("Uncertain")\\
(from now on, we will forget about "$2\times$" form the definition, for simplicity)
\item (of valency 1):\\
"(X unique)" corresponds to the collection of all conglomerate--semisingletones.\\
"(X exists)" corresponds to the collection of all nonempty conglomerates.\\
"(X (well--)defined)" corresponds to the collection of all nonempty conglomerate--semisingletones, i.e. to the collection of all conglomerate--singletones\\
"empty" \ corresponds to the collection of empty conglomerates, i.e. just $\{\emptyclass\}$.
\item (of valency 2):\\
"$(x\in y)$" corresponds to collection of all ordered pairs of nonempty conglomerates, where each class from first conglomerate is in each class of the second conglomerate.\\
"$(x\notin y)$" corresponds to collection of all ordered pairs of nonempty conglomerates, where each class from first conglomerate is not in any class of the second conglomerate.
\end{itemize}
\begin{definition}(Theorem)
Theorem is a predicate which has the truth value "True".\\
\end{definition}
%$\tau t$
%Therefore.
%This is something unresolved, we need some time...


% and define the Uncertain value, like:\\
%Uncertain = Non--existent OR Non--constant\\
%\,\quad\quad = Absent OR Unstable\\
%\,\quad\quad = Void OR Varying\\
%But then, how connectives will work? They will be "contagious" take "the worst case"? %Not necessary: two predicates can stabilize each other.\\
Idea 2:\\

Actually, variables are ordered: $x_1,x_2,x_3,\dots$.
So we can consider the following correspondence evaluation of the formula $x_2\in x_4$ via assignment function $w$, with
\[
(1\times)\Un\times \{w(x_2)\}\times\Un\times\{w(x_4)\}\times\Un\times\Un\times\Un\dots,
\]
where "infinite product" is abuse of notation, which means a collection of functions $f$ from $\mathbb{N}_0$ to $\power\Un$, such that $f(0)=1$. %\subseteq

\iffalse
\symbol{200}\symbol{201}\symbol{202}\symbol{203}
\symbol{204}\symbol{205}\symbol{206}\symbol{207}
\symbol{208}\symbol{209}\symbol{210}\symbol{211}
\symbol{212}\symbol{213}\symbol{214}\symbol{215}
\symbol{216}\symbol{217}\symbol{218}\symbol{219}\\
\symbol{220}\symbol{221}\symbol{222}\symbol{223}
\symbol{224}\symbol{225}\symbol{226}\symbol{227}
\symbol{228}\symbol{229}\symbol{230}\symbol{231}
\symbol{232}\symbol{233}\symbol{234}\symbol{235}
\symbol{236}\symbol{237}\symbol{238}\symbol{239}\\
\symbol{240}\symbol{241}\symbol{242}\symbol{243}
\symbol{244}\symbol{245}\symbol{246}\symbol{247}
\symbol{248}\symbol{249}\symbol{250}\symbol{251}
\symbol{252}\symbol{253}\symbol{254}\symbol{255}
\fi

%\symbol{253A}
\end{comment}

\ifdefined\qwertyflag
\undef\qwertyflag
\end{document}
\fi
